% Imporditud: tools/import_eksperimendid.py
% Allikas: Eesti füüsikaolümpiaadi eksperimendi ülesannete
%   kogu (Rael Kalda, 2023), ülesanne Ü124, lahendus L124.
\setAuthor{Kaarel Kivisalu}
\setRound{piirkonnavoor}
\setYear{2022}
\setNumber{G E1}
\setDifficulty{3}
\setTopic{dünaamika}
\setVahendid{1 tavaline spagett (spageti purunemisel on võimalik juurde saada vähemalt 9 tervet spagetti), millimeeterpaber}
\setPunktid{10}
\setAllikas{Eksperimendi kogu Ü124, lahendus lk 98}
\setLahendusseis{toores}

\prob{Spagett}
Hinnake, kui pikk peaks minimaalselt olema spagett, et selle otspunktid kokku pannes moodustuks ringjoon.

\hint

\solu
Lahendus 1: Painutame spagetikõrt võttes mõlemast otsast kahe näpuga kinni ja suurendades otste (täpsemalt: otstest tõmmatud puutujate) vahelist nurka $\alpha$ kuni spagetikõrre murdumiseni ning teeme kindlaks vastava maksimaalse nurga $\alpha$max. Suure spagetiringi raadius on siis R = l/$\alpha$max, kus l on kõrre pikkus.

Märkus: Nurga mõõtmine on eksperimentaalselt suhteliselt keeruline kasutades ainult millimeeterpaberit, kuid trigonomeetriat kasutades siiski võimalik.

Lahendus 2:

Painutame spagetikõrt analoogselt lahendusega 1. Seekord mõõdame vabalt sõr-

mede vahelise spagetikõrre pikkuse l ning kui kaugele spageti keskpunkt maksi-

maalselt jõuab kinnihoidmiskohtasid ühendavast sirgest a. Siis saab kõversuraa-

diuse leida võrrandist 1$-$ a = cos L . Seda võrrandit tuleb lahendada numbriliselt, R 2R kuna vähemalt lihtsat analüütilist lahendust pole.

HIndamine:

Spageti painutamine nii, et avaldatakse näppudega ainult jõumomenti (mitte jõudu) spageti painutamiseks. Sellisel juhul moodustub spagetist ringjoone kaar. (3 p)

Märkus: Rakendades jõudu ei teki spagetis ringjoone kaart. Siiski sobivalt jõudu rakendades on võimalik kirjeldada spageti kuju (sinusoid, kuupparabool vms sõltuvalt jõust ja jõumomendist) ning leida Youngi moodul, mille kaudu saab avaldada ringikujulise spageti raadiuse. Koos piisava matemaatilise kirjeldusega on võimalik ka sellisel juhul saada täispunktid.

Oma valitud meetodi jaoks leitud võrrand raadiuse leidmiseks. (3 p) Tehtud vähemalt 3 mõõtmist (punktide saamiseks vajalik andmepunktide ole-

masolu). (2 p) Iga mõõtmise jaoks välja arvutatud pika spageti raadius. (1 p) Lõppvastus mõistlikus vahemikus keskmisena üksikute mõõtmiste raadiustest.

Ligikaudu 850 mm, see võib varieeruda arvestatavalt kuna spagettide tootjad on erinevad ning esineb ka spagetikõrte vaheline variatsioon. Vastused võiksid olla vahemikus 700 mm kuni 1000 mm (ideaalis peaks iga spagettide tootja jaoks tegema kontrollmõõtmised ja selle järgi määrama sobiliku vahemiku). (1 p)
\probend
